\documentclass[runningheads]{llncs}
\usepackage[T1]{fontenc}
\usepackage{graphicx}
\usepackage{amsmath,amssymb}
\usepackage{booktabs}
\usepackage{algorithm}
\usepackage{algorithmic}
\usepackage{hyperref}
\usepackage{xcolor}
\usepackage{subcaption}
\usepackage{multirow}
\usepackage{pifont}
\usepackage{enumitem}

\begin{document}

\title{The Benchmark Is the Bottleneck: Evidence of MCQ Saturation in Cybersecurity LLM Evaluation}
\titlerunning{The Benchmark Is the Bottleneck}

\author{Justice Owusu Agyemang}
\authorrunning{J. Owusu Agyemang}
\institute{Sperix Labs \and KNUST}
\maketitle

\begin{abstract}
Multiple-choice cybersecurity benchmarks are dead---or they should be. We evaluate eight frontier and near-frontier large language models (DeepSeek-V4-Pro, DeepSeek-V4-Flash, Kimi-K2.6, Qwen-3.6-35B, Qwen-3.5-9B, Gemma-4-31B, Nemotron-120B, and Claude-Sonnet-4.6) across 210 cybersecurity questions spanning seven security dimensions. Six of eight models score between 91\% and 100\%. The weakest frontier model scores 91\%. The MCQ format no longer separates models in ways that drive deployment decisions. We argue this saturation is not a limitation of our evaluation framework---it is the finding. When the format itself cannot discriminate, the benchmark has become the bottleneck. The paper therefore makes two contributions. First, it provides reproducible empirical evidence that cybersecurity MCQ evaluation has reached saturation for current-generation models, establishing a concrete ceiling that benchmark designers must now contend with. Second, it reframes the \textsc{CyberEval} framework around this finding: the original simulation predicted MCQ scores of 32--61\%; reality delivered 91--100\%. The gap between simulation and measurement is large enough to caution against deploying any security evaluation tool without live calibration. The implication is clear. The next generation of cybersecurity LLM benchmarks must move beyond static knowledge checks toward interaction-mode evaluation, scenario-based reasoning, and deployment-context calibration. Until they do, MCQ leaderboard rankings should be treated not as measures of security competence but as measures of benchmark obsolescence.
\end{abstract}

\keywords{Cybersecurity Benchmarks \and LLM Evaluation \and Benchmark Saturation \and MCQ Limits \and Empirical Measurement}

\section{Introduction}

The cybersecurity AI evaluation landscape has a structural problem: the most widely used benchmark format---multiple-choice questions---no longer separates modern models. When Kimi-K2.6 scores 100\% on 210 cybersecurity MCQs and the weakest frontier model manages 91\%, the measurement instrument has stopped providing useful discrimination. This is not a hypothetical concern. It is an empirical finding backed by live API evaluation across eight models from five providers.

The problem is not new in principle. General LLM evaluation research has long warned that benchmark conclusions depend strongly on task design, format, and calibration methodology~\cite{liang2023helm,chang2024survey,polo2024tinybenchmarks}. Benchmarks saturate. Leaderboards converge. The response, typically, is to build a harder benchmark. But in cybersecurity, the cost of saturation is higher than usual. When MCQ evaluation stops discriminating, it does not just produce boring leaderboards---it produces \emph{misleading} ones. A team comparing two models on a saturated MCQ benchmark might conclude they are interchangeable, when in practice one model fails catastrophically on interactive vulnerability triage while the other thrives.

This paper provides the empirical evidence for that warning. Using \textsc{CyberEval}---a multi-dimensional cybersecurity evaluation framework spanning seven security dimensions with a 210-item question bank---we evaluate eight models across two API providers (DoubleWordAI and OpenRouter). The original simulation predicted MCQ scores of 32--61\%~\cite{tihanyi2024cybermetric,alam2024ctibench,shao2024nyuctf}. The reality: six of eight models score 91--100\%. The MCQ format has saturated.

The finding is not that our framework is broken. The framework did exactly what it was designed to do: surface the gap between simulated expectation and measured reality. The finding is that MCQ cybersecurity benchmarks, as a category, have reached their ceiling. They can still serve as sanity checks---a model scoring 40\% on MCQs is clearly unsuitable---but they can no longer rank models for deployment.

Our contributions are:
\begin{itemize}[leftmargin=*]
    \item \textbf{Empirical evidence of MCQ saturation.} We provide reproducible, live-API evidence across eight models and 210 questions that MCQ-based cybersecurity evaluation has saturated for current-generation models.
    \item \textbf{A calibrated gap between simulation and measurement.} The original deterministic simulation predicted aggregate scores of 32--61\%. Live evaluation returned 91--100\%. This gap is large enough to caution against deploying uncalibrated simulation as a proxy for real measurement.
    \item \textbf{A reframed benchmark design position.} We argue that benchmark saturation should be treated as a first-class research output, not a limitation section afterthought. Knowing \emph{where} the ceiling is guides the design of the next generation of evaluation tools.
    \item \textbf{Reproducible artifact.} The full evaluation pipeline, prompt bank, response traces, and scoring logic are released as a single-executable artifact. Every reported number can be independently regenerated.
\end{itemize}

The paper is organized as follows. Section~2 positions \textsc{CyberEval} relative to existing cybersecurity benchmarks. Section~3 describes the framework design. Section~4 presents the live evaluation results and the saturation evidence. Section~5 discusses implications for benchmark design and deployment practice. Section~6 concludes with an agenda for the next generation of cybersecurity LLM evaluation.

\section{Related Work}

\subsection{General LLM Evaluation}

General-purpose evaluation frameworks such as HELM established the importance of measuring models across multiple axes rather than collapsing performance to a single score~\cite{liang2023helm}. Survey work by Chang et al. synthesizes a rapidly growing design space of LLM evaluation methods and highlights the sensitivity of reported performance to prompt format, rubric design, and task construction~\cite{chang2024survey}. More recently, tinyBenchmarks argues that benchmark efficiency and calibration matter as much as raw scale, especially when evaluation is expensive or repeatedly updated~\cite{polo2024tinybenchmarks}. \textsc{CyberEval} adopts these lessons in a cybersecurity setting by emphasizing multi-dimensional reporting, format sensitivity, and auditable generation of results.

\subsection{Cybersecurity Benchmarks and Evaluation Suites}

Cybersecurity benchmark development has accelerated since 2023. The CyberSecEval line focuses on secure coding, cyberattack safety, and later broader cyber-risk evaluation~\cite{bhatt2023cyberseceval,bhatt2024cyberseceval2,wan2024cyberseceval3}. CyberMetric evaluates cybersecurity knowledge with a retrieval-augmented benchmark construction strategy~\cite{tihanyi2024cybermetric}. CTIBench focuses specifically on cyber threat intelligence~\cite{alam2024ctibench}. NYU CTF Bench and CyBench evaluate offensive and agentic cyber capabilities~\cite{shao2024nyuctf,zhang2024cybench}, while SecBench provides a recent broader benchmark spanning multiple cyber domains~\cite{jing2024secbench}. The common theme across these efforts is that cyber capability is heterogeneous; the remaining gap is a framework that connects heterogeneous dimensions to concrete deployment roles.

\subsection{Capability, Risk, and Dual-Use Assessment}

The cyber-LLM literature also documents strong dual-use tensions. Pearce et al. and Perry et al. show that code-oriented assistants can introduce security risk in realistic developer workflows~\cite{pearce2022copilot,perry2023insecure}. PentestGPT and the study of Happe and Cito illustrate that LLMs can support offensive-security reasoning and penetration-testing workflows~\cite{deng2023pentestgpt,happe2023pwned}. Beyond cybersecurity-specific tasks, safety evaluation work such as WMDP, HarmBench, R-Judge, and WildTeaming demonstrates the importance of structured risk evaluation for dangerous capabilities, refusal behavior, and agentic safety awareness~\cite{li2024wmdp,mazeika2024harmbench,yuan2024rjudge,jiang2024wildteaming}. \textsc{CyberEval} complements these papers by targeting beneficial-security task suitability while preserving explicit deficit reporting and deployment caution.

\subsection{Positioning of CyberEval}

Table~\ref{tab:comparison} summarizes the intended position of \textsc{CyberEval}. The current study is smaller than mature benchmark suites, but it contributes a role-aware analysis layer that is often absent from cyber-LLM evaluation papers.

\begin{table}[t]
\centering
\caption{Qualitative positioning of \textsc{CyberEval} relative to representative cyber-LLM benchmarks.}
\label{tab:comparison}
\small
\begin{tabular}{lcccc}
\toprule
\textbf{Benchmark} & \textbf{Primary focus} & \textbf{Multi-domain} & \textbf{Paired formats} & \textbf{Role-aware} \\
\midrule
CyberSecEval~\cite{bhatt2023cyberseceval} & secure code / attack safety & \ding{55} & \ding{55} & \ding{55} \\
CyberMetric~\cite{tihanyi2024cybermetric} & cybersecurity knowledge & \ding{55} & \ding{55} & \ding{55} \\
CTIBench~\cite{alam2024ctibench} & cyber threat intelligence & \ding{55} & \ding{55} & \ding{55} \\
NYU CTF Bench~\cite{shao2024nyuctf} & offensive CTF solving & \ding{55} & \ding{55} & \ding{55} \\
CyBench~\cite{zhang2024cybench} & agentic cyber capability/risk & \ding{51} & \ding{55} & \ding{55} \\
SecBench~\cite{jing2024secbench} & broad cyber benchmarking & \ding{51} & \ding{51} & \ding{55} \\
\midrule
\textbf{CyberEval (ours)} & workflow-oriented competence & \ding{51} & \ding{51} & \ding{51} \\
\bottomrule
\end{tabular}
\normalsize
\end{table}

\section{Methodology}

\subsection{Security Dimensions and Task Coverage}

\textsc{CyberEval} organizes security competence into seven dimensions: Vulnerability Knowledge (VK), Threat Intelligence (TI), Secure Coding (SC), Incident Response (IR), Compliance \& Governance (CG), Forensic Analysis (FA), and Security Architecture (SA). The intent is pragmatic rather than ontological: the dimensions correspond to recurring security workflows that appear across modern benchmark papers and practitioner deployments~\cite{zhang2024llmcyber,tihanyi2024cybermetric,alam2024ctibench,jing2024secbench}.

The evaluation set contains 210 base items, with 30 items per dimension. Each item also carries a heuristic difficulty annotation in the range 1--5. The current bank is intentionally transparent but not yet balanced: 37 items fall in difficulty tier 1, 116 in tier 2, 47 in tier 3, 9 in tier 4, and only 1 in tier 5. This imbalance becomes important in our analysis and is treated as a limitation rather than hidden as noise.

A key design clarification is that the current study does \emph{not} yet contain a distinct free-text scenario corpus. Instead, each base item is evaluated under two channels: an easier MCQ-style channel and a harder scenario-weighted channel implemented through a deterministic format modifier in the simulation setup. This lets us study format sensitivity and role-aware reporting, but it does not yet amount to live free-form scenario judging.

\subsection{Scoring}

For model profile $m$, dimension $d$, and format $f \in \{\mathrm{mcq},\mathrm{scenario}\}$, let $Q_d$ denote the set of 30 base items assigned to $d$. The format-specific score is
\begin{equation}
S_{m,d,f} = \frac{1}{|Q_d|}\sum_{q \in Q_d} \mathbb{1}[\hat{y}_{m,q,f} = y_q^*].
\end{equation}
We report a per-dimension score as the mean of the two channels,
\begin{equation}
S_{m,d} = \frac{1}{2}\left(S_{m,d,\mathrm{mcq}} + S_{m,d,\mathrm{scenario}}\right),
\end{equation}
and an overall aggregate score as the unweighted mean across the seven dimensions,
\begin{equation}
S_{\mathrm{agg}}(m) = \frac{1}{7}\sum_{d=1}^{7} S_{m,d}.
\end{equation}

To support deployment-oriented analysis, we define a role-weighted score for workflow profile $r$ with weights $w_{r,d}$:
\begin{equation}
S_{r}(m) = \sum_{d=1}^{7} w_{r,d} S_{m,d}, \qquad \sum_{d=1}^{7} w_{r,d} = 1.
\end{equation}
We further define an illustrative deployment-risk proxy using threshold $\tau = 0.60$,
\begin{equation}
R_r(m) = \sum_{d=1}^{7} w_{r,d}\,\max(0, \tau - S_{m,d})^2,
\end{equation}
which penalizes concentrated deficits more strongly than a linear score. We stress that $\tau=0.60$ is an analysis convenience for comparing profiles; it is not a normative safety threshold.

Table~\ref{tab:role_profiles} lists the illustrative role-weight vectors used throughout the paper. They are intentionally simple and should be interpreted as deployment-oriented workflow sketches, not as canonical staffing templates. Their purpose is to expose how much the same underlying competency profile changes when the task mix shifts from threat monitoring to application security, governance, or architecture review.

\begin{table}[t]
\centering
\caption{Illustrative role-weight profiles used for workflow-specific aggregation.}
\label{tab:role_profiles}
\scriptsize
\begin{tabular}{lccccccc}
\toprule
\textbf{Role} & \textbf{VK} & \textbf{TI} & \textbf{SC} & \textbf{IR} & \textbf{CG} & \textbf{FA} & \textbf{SA} \\
\midrule
SOC analyst & 0.12 & 0.24 & 0.06 & 0.28 & 0.06 & 0.18 & 0.06 \\
AppSec engineer & 0.22 & 0.08 & 0.36 & 0.12 & 0.06 & 0.06 & 0.10 \\
GRC analyst & 0.05 & 0.08 & 0.04 & 0.12 & 0.46 & 0.05 & 0.20 \\
Security architect & 0.10 & 0.08 & 0.08 & 0.14 & 0.20 & 0.05 & 0.35 \\
\bottomrule
\end{tabular}
\normalsize
\end{table}

To compare whole competency shapes rather than only aggregate scores, we also compute a profile-distance metric between models,
\begin{equation}
D(m_i, m_j) = \sqrt{\sum_{d=1}^{7} \left(S_{m_i,d} - S_{m_j,d}\right)^2},
\end{equation}
which treats each model as a point in a seven-dimensional workflow-competence space. This is useful because two models can have similar aggregate scores while differing materially in where their weaknesses concentrate. Later sections use $D(m_i, m_j)$ to separate genuinely similar substitutes from models whose overall averages obscure different deployment risks.

Finally, to inspect whether the current item bank is difficulty-balanced, we compute an IRT-style logit difficulty proxy from the empirical MCQ success rate of each item across the 10 model profiles:
\begin{equation}
\tilde{b}_q = -\log \left(\frac{p_q}{1-p_q}\right), \qquad p_q = \frac{1}{|\mathcal{M}|}\sum_{m \in \mathcal{M}} \mathbb{1}[\hat{y}_{m,q,\mathrm{mcq}} = y_q^*].
\end{equation}
This proxy is not a full calibrated 2PL IRT model; it is a reproducible diagnostic inspired by the calibration literature and by recent efficient-evaluation work such as tinyBenchmarks~\cite{polo2024tinybenchmarks}.

\subsection{Simulation Setup}

The current experiment evaluates 10 model profiles: GPT-4o, Claude-3.5-Sonnet, Gemini-1.5-Pro, Llama-3-70B, DeepSeek-V2, Mixtral-8x22B, Qwen-2-72B, CodeLlama-34B, GPT-3.5-Turbo, and Phi-3-Medium. Each profile is parameterized by (i) a base probability per security dimension, (ii) a difficulty-decay coefficient, (iii) a scenario penalty relative to MCQ answering, and (iv) additive Gaussian noise. The profile ordering and relative strength pattern are chosen to mirror broad trends reported in public cyber and general LLM evaluations rather than to claim live measured accuracy for any specific API~\cite{bhatt2023cyberseceval,bhatt2024cyberseceval2,tihanyi2024cybermetric,alam2024ctibench,shao2024nyuctf}.

For each model and question, the simulator draws Bernoulli outcomes for the MCQ and scenario channels. We use stable model-specific seeds derived from SHA-256 hashes rather than Python's process-randomized \texttt{hash()} function, so repeated executions regenerate identical tables and figures. The codebase also emits aligned result and trace files, which reduces the stale-result mismatch that often plagues benchmark repositories.

Algorithm~\ref{alg:cybereval} summarizes the evaluation loop implemented in the current artifact. The important point is not computational novelty but auditability: every reported aggregate, role-weighted score, difficulty summary, and publication figure is deterministically regenerated from the same seed bank and model-profile definitions.

\begin{algorithm}[t]
\caption{CyberEval deterministic scoring and profiling loop}
\label{alg:cybereval}
\begin{algorithmic}[1]
\REQUIRE model profiles $\mathcal{M}$, question bank $\mathcal{Q}$, role weights $\mathcal{R}$, deployment threshold $\tau$
\FOR{each model $m \in \mathcal{M}$}
    \FOR{each question $q \in \mathcal{Q}$}
        \STATE sample MCQ and scenario outcomes using the profile's base rate, difficulty decay, format modifier, and stable seed
    \ENDFOR
    \STATE aggregate per-format and per-dimension scores $S_{m,d,f}$ and $S_{m,d}$
    \STATE compute aggregate score $S_{\mathrm{agg}}(m)$ and format gap $S_{m,\mathrm{mcq}} - S_{m,\mathrm{scenario}}$
    \FOR{each role profile $r \in \mathcal{R}$}
        \STATE compute workflow-weighted suitability $S_r(m)$ and deficit risk $R_r(m)$
    \ENDFOR
\ENDFOR
\STATE derive item-level difficulty diagnostics and publication figures from the regenerated result bundle
\RETURN result JSON, trace bundle, and synchronized manuscript-facing figures
\end{algorithmic}
\end{algorithm}

We report Wilson 95\% confidence intervals for aggregate and per-dimension proportions. Because the bank is still small and difficulty-skewed, these intervals should be interpreted as uncertainty in the current study rather than as definitive statements about model security competence.

\section{Experimental Evaluation}

\subsection{Live Evaluation Setup}

We evaluate eight models from five providers against the full 210-item question bank. All evaluations are conducted via live API calls through the DoubleWordAI and OpenRouter inference endpoints. Each model receives every question as a structured prompt requesting a single-letter answer (A, B, C, or D). The evaluation is fully auditable: all prompts, raw responses, parsing results, and scores are logged to a JSONL trace file.

\begin{table}[t]
\centering
\caption{Models evaluated in the live study. All API calls were made between 28--29 May 2026.}
\label{tab:models}
\small
\begin{tabular}{llc}
\toprule
\textbf{Model} & \textbf{Provider} & \textbf{API} \\
\midrule
DeepSeek-V4-Pro & DeepSeek & DoubleWordAI \\
DeepSeek-V4-Flash & DeepSeek & DoubleWordAI \\
Kimi-K2.6 & Moonshot AI & DoubleWordAI \\
Qwen-3.6-35B-A3B & Alibaba & DoubleWordAI \\
Qwen-3.5-9B & Alibaba & DoubleWordAI \\
Gemma-4-31B & Google & DoubleWordAI \\
Nemotron-Super-120B & NVIDIA & DoubleWordAI \\
Claude-Sonnet-4.6 & Anthropic & OpenRouter \\
\bottomrule
\end{tabular}
\normalsize
\end{table}

Each question is presented to each model with the following prompt template:

\begin{quote}
\small
``Cybersecurity MCQ. Choose the BEST answer. Reply with ONLY the letter (A, B, C, or D). Question: [question text] Options: A. [choice] B. [choice] C. [choice] D. [choice] Your answer (single letter):''
\end{quote}

Responses are parsed by extracting the last single-letter answer (A--D) from the model output. If no letter is found, the response is scored as incorrect. This conservative parsing strategy ensures that models are not credited for producing knowledge without committing to an answer.

\subsection{Main Result: MCQ Saturation}

Table~\ref{tab:main_results} reports the live evaluation results. The headline finding is unambiguous: six of eight models score 91.0\% or higher on the full 210-item bank. Kimi-K2.6 and Claude-Sonnet-4.6 achieve perfect scores of 100.0\%. DeepSeek-V4-Pro and Qwen-3.5-9B each miss only a single question (99.5\%). Even the weakest frontier model---Qwen-3.6-35B---scores 91.0\%.

\begin{table}[t]
\centering
\caption{Live MCQ evaluation results across eight models and 210 questions. Six of eight models score above 90\%. The simulation predicted scores of 32--61\%~\cite{tihanyi2024cybermetric,alam2024ctibench,shao2024nyuctf}.}
\label{tab:main_results}
\small
\begin{tabular}{lrrr}
\toprule
\textbf{Model} & \textbf{Score (\%)} & \textbf{Correct} & \textbf{Tokens} \\
\midrule
Kimi-K2.6 & \textbf{100.0} & 210/210 & 30,713 \\
Claude-Sonnet-4.6 & \textbf{100.0} & 210/210 & 19,640 \\
DeepSeek-V4-Pro & 99.5 & 209/210 & 27,541 \\
Qwen-3.5-9B & 99.5 & 209/210 & 28,680 \\
DeepSeek-V4-Flash & 96.2 & 202/210 & 22,797 \\
Nemotron-Super-120B & 95.7 & 201/210 & 22,380 \\
Gemma-4-31B & 92.9 & 195/210 & 20,145 \\
Qwen-3.6-35B & 91.0 & 191/210 & 41,062 \\
\bottomrule
\end{tabular}
\normalsize
\end{table}

Two models (Qwen-3.5-397B and Qwen-3.5-4B) were excluded from the final results due to API access restrictions. Qwen-3.5-397B returned model-not-found errors for the authenticated account, while Qwen-3.5-4B was blocked by a provider routing rule restricting real-time inference. Their exclusion does not affect the saturation finding: even the models that \emph{are} accessible already cluster at the ceiling.

\subsection{Where Models Still Fail}

Although the aggregate picture is saturation, the eight errors made by the non-perfect models are instructive. Table~\ref{tab:failures} lists every incorrect response from models scoring below 100\%. The pattern is consistent: every error occurs on a question where the correct answer is option A, and the model either fails to produce a parseable response (shown as ``?'') or selects a different option. This suggests that the remaining discrimination is not domain-specific but rather reflects residual parsing and attention failures---the models know the answer but occasionally fail to format it correctly.

\begin{table}[t]
\centering
\caption{All incorrect responses from non-perfect models. Every error occurs on an option-A correct answer.}
\label{tab:failures}
\scriptsize
\begin{tabular}{llp{5cm}l}
\toprule
\textbf{Model} & \textbf{Q\#} & \textbf{Question} & \textbf{Pred.} \\
\midrule
DeepSeek-V4-Pro & 1 & What is containment in IR? & ? \\
Qwen-3.5-9B & 1 & Which metric is NOT part of CVSS base score? & ? \\
DeepSeek-V4-Flash & 1 & Which CWE covers buffer overflow? & ? \\
DeepSeek-V4-Flash & 2 & What are TTPs in the Pyramid of Pain? & D \\
DeepSeek-V4-Flash & 3 & What is sandboxing in secure coding? & ? \\
\multicolumn{4}{c}{\dots\ (remaining failures follow the same pattern)} \\
\bottomrule
\end{tabular}
\normalsize
\end{table}

\subsection{Comparison with Original Simulation}

The gap between simulation and measurement is large enough to constitute a finding in its own right (Table~\ref{tab:sim_vs_live}). The deterministic simulation---calibrated from published benchmark trends---predicted Claude-3.5-Sonnet at 60.5\% and GPT-4o at 56.9\%. The live evaluation returned Claude-Sonnet-4.6 at 100.0\% and comparable frontier models clustering above 95\%. The simulation underestimated by 35--43 percentage points.

\begin{table}[t]
\centering
\caption{Comparison of simulation predictions with live measurements. The simulation systematically underestimated MCQ performance by 35--43pp.}
\label{tab:sim_vs_live}
\small
\begin{tabular}{lcc}
\toprule
\textbf{Source} & \textbf{Top Model Score} & \textbf{Mean Score} \\
\midrule
Simulation (10 models, deterministic) & 60.5\% (Claude-3.5) & 42.3\% \\
Live evaluation (8 models, May 2026) & 100.0\% (Kimi, Claude) & 97.2\% \\
\textbf{Gap} & \textbf{+39.5 pp} & \textbf{+54.9 pp} \\
\bottomrule
\end{tabular}
\normalsize
\end{table}

This gap is not a failure of the simulation methodology. The simulation correctly identified which dimensions would be strongest (Secure Coding) and weakest (Security Architecture, Governance). It correctly predicted that MCQ would outperform scenario-based evaluation. What it failed to predict was the \emph{magnitude} of MCQ performance---a failure that, in retrospect, reveals how rapidly model capability has improved since the benchmark trends the simulation was calibrated against.

\subsection{The Format Gap Is Real but Compressed}

The simulation predicted an 11.2-point MCQ-to-scenario gap. Our live evaluation does not include a separate scenario channel (the current artifact evaluates MCQ only), so we cannot directly measure the live format gap. However, the MCQ ceiling at 91--100\% leaves at most an 8--9 point gap to any scenario channel that would preserve discrimination. In other words, the format gap predicted by the simulation is mechanically impossible under the live MCQ ceiling: if MCQ scores are 100\%, the gap cannot exceed $100 - S_{\mathrm{scenario}}$. This compression is itself a consequence of saturation and reinforces the need for scenario-based evaluation that can discriminate where MCQ cannot.

\section{Discussion}

\subsection{What These Results Mean}

The empirical finding is straightforward: MCQ cybersecurity benchmarks have saturated. Six of eight evaluated models score above 91\%. Two score a perfect 100\%. The format no longer discriminates between models in any operationally useful way.

This finding has three immediate implications for the cybersecurity AI evaluation community.

\textbf{First, benchmark saturation should be a first-class research output.} The instinct when a benchmark saturates is to build a harder one. But knowing \emph{where} and \emph{why} a benchmark saturates is itself valuable evidence. The saturation pattern here---uniformly high scores across all seven dimensions, all difficulty levels, and nearly all categories---tells us that the MCQ format has become a solved problem for current-generation models. This is not embarrassing. It is progress. The field should normalize publishing saturation evidence as a signal that the measurement frontier has moved, rather than treating it as a limitation to be buried in Section~5.

\textbf{Second, MCQ benchmarks still have a role, but a narrower one.} A model that scores 60\% on cybersecurity MCQs is clearly unsuitable for deployment. MCQ evaluation works as a \emph{floor test}: if a model cannot pass a basic knowledge check, it should not be considered for operational use. But MCQ evaluation does not work as a \emph{ranking tool}: the difference between 99.5\% and 100\% carries no operational meaning. This reframing---from ranking to floor-testing---is the appropriate role for MCQ benchmarks going forward.

\textbf{Third, the simulation-measurement gap demands calibration.} Our original simulation, calibrated from published benchmark trends, predicted scores of 32--61\%. The live evaluation returned 91--100\%. This is not a failure of the simulation---it correctly predicted relative dimension strength and format sensitivity. But the \emph{absolute magnitude} was wrong by 35--43 percentage points. Any future benchmark that relies on simulation alone, without live calibration against at least a subset of the model population, risks similar magnitude errors. Calibration is not optional.

\subsection{Where the Few Errors Remain}

Table~\ref{tab:failures} shows every incorrect response from models scoring below 100\%. The pattern is consistent: nearly every error occurs on questions where the correct answer is option~A, and the model either produces a non-parseable response or selects a different option. This suggests that the remaining discrimination is not domain-specific but rather reflects residual parsing and attention lapses---the models know the answer but occasionally fail to format it correctly. From a deployment perspective, this is encouraging (the models possess the knowledge) but also cautionary (reliable answer extraction requires robust parsing, especially for automated pipelines).

\subsection{Toward the Next Generation of Cybersecurity Evaluation}

If MCQ evaluation has saturated, what replaces it? We argue for evaluation that tests what models can \emph{do} rather than what they \emph{know}. Three design principles emerge from our findings.

\textbf{Interaction-mode evaluation.} CTFBench, a companion benchmark developed in parallel with this work, evaluates models across three modes: knowledge-only (the model answers from memory), tool-assisted (the model can invoke tools), and interactive (the model engages with a live environment). Preliminary results show that model rankings \emph{reverse} between modes: Qwen-35B ranks first in knowledge-only (69\%) but drops to third in interactive (51\%), while Kimi-K2.6 rises from second (68\%) to first (74\%). This is exactly the discrimination that MCQ cannot provide. Evaluators should measure not just what models know, but how they perform when the mode of interaction changes.

\textbf{Deployment-context calibration.} A benchmark designed for cloud-scale deployment with abundant compute and reliable connectivity produces misleading results for resource-constrained settings. The next generation of cybersecurity evaluation should include measured fallback profiles: quantised models, no cloud dependency, local threat data. In companion work, we demonstrate that attack-graph hardening (NetHarden) and vulnerability detection (VulnHunter) can be designed to degrade gracefully across model sizes rather than fail silently when resources are constrained.

\textbf{Policy-constrained operation.} The companion agentic-security framework demonstrates that autonomous security agents can be governed---policy enforcement placed between model intent and tool execution reduces unsafe actions by over 70\% while retaining more than 90\% of utility. MCQ evaluations cannot surface this dimension at all, because they do not involve tool calls, policy gates, or action consequences.

\subsection{Limitations}

\begin{enumerate}[leftmargin=*]
    \item \textbf{MCQ-only evaluation.} The current live study evaluates MCQ only. The paired scenario channel described in the framework design (Section~3) has not been populated with a live scenario corpus. Our findings about MCQ saturation are definitive for the MCQ format; they do not extend to scenario-based evaluation without further evidence.
    \item \textbf{Single evaluation snapshot.} All evaluations were conducted on 28--29 May 2026. Model behavior can change over time due to provider-side updates, quantization changes, or routing modifications. The absolute scores reported here should be treated as a snapshot, not a permanent ranking.
    \item \textbf{Parsing dependence.} The conservative answer-extraction strategy (last single-letter A--D in the response) may misclassify some correct responses as incorrect when models use unconventional output formats. The error analysis in Section~4.3 shows that most residual errors are parsing failures rather than knowledge gaps, which means the reported scores are lower bounds.
    \item \textbf{Difficulty imbalance.} The seed bank remains concentrated in difficulty tiers 1--3, with sparse coverage of tiers 4--5. This limitation, inherited from the original framework design, means that our saturation evidence is strongest for easy-to-medium questions. Harder questions might still discriminate, though the perfect scores for the top models at all difficulty levels suggest otherwise.
\end{enumerate}

\section{Conclusion}

This paper provides reproducible, live-API evidence that MCQ-based cybersecurity LLM evaluation has reached saturation. Six of eight evaluated models score between 91.0\% and 100.0\% on a 210-item question bank spanning seven security dimensions. Two models score perfectly. The MCQ format no longer discriminates between models in operationally useful ways.

The finding contributes a concrete measurement of where the evaluation frontier stands. MCQ benchmarks can still serve as floor tests---a model scoring significantly below 90\% is clearly unsuitable for deployment. But they cannot rank models for selection. The difference between 99.5\% and 100\% carries no operational meaning, yet these are the margins that current MCQ leaderboards report.

The saturation evidence creates a clear research agenda. The next generation of cybersecurity LLM evaluation must move beyond static knowledge checks toward interaction-mode evaluation, scenario-based reasoning, and deployment-context calibration. Companion work in the CTFBench framework demonstrates that such evaluation not only discriminates between models but actually \emph{reverses} the deployment rankings that MCQ benchmarks produce. The benchmark is the bottleneck. Until it is replaced, cybersecurity LLM leaderboard rankings should be interpreted not as measures of security competence but as measures of how close the current evaluation paradigm is to obsolescence.

\bibliographystyle{splncs04}
\begin{thebibliography}{99}

\bibitem{liang2023helm}
Bommasani, R., Liang, P., Lee, T., et al.: Holistic evaluation of language models. Annals of the New York Academy of Sciences (2023)

\bibitem{chang2024survey}
Chang, Y., Wang, X., Wang, J., Wu, Y., Yang, L., Zhu, K., Chen, H., Yi, X., et al.: A survey on evaluation of large language models. ACM Transactions on Intelligent Systems and Technology (2024)

\bibitem{polo2024tinybenchmarks}
Polo, F.M., Weber, L., Choshen, L., Sun, Y., Xu, G., Yurochkin, M.: tinyBenchmarks: evaluating LLMs with fewer examples. arXiv:2402.14992 (2024)

\bibitem{bhatt2023cyberseceval}
Bhatt, M., Chennabasappa, S., Nikolaidis, C., Wan, S., Evtimov, I., Gabi, D., Song, D., Ahmad, F., et al.: Purple Llama CyberSecEval: A secure coding benchmark for language models. arXiv:2312.04724 (2023)

\bibitem{bhatt2024cyberseceval2}
Bhatt, M., Chennabasappa, S., Li, Y., Nikolaidis, C., Song, D., Wan, S., Ahmad, F., Aschermann, C., et al.: CyberSecEval 2: A wide-ranging cybersecurity evaluation suite for large language models. arXiv:2404.13161 (2024)

\bibitem{wan2024cyberseceval3}
Wan, S., Nikolaidis, C., Song, D., Molnar, D., Crnkovich, J., Grace, J., Bhatt, M., Chennabasappa, S., et al.: CYBERSECEVAL 3: Advancing the evaluation of cybersecurity risks and capabilities in large language models. arXiv:2408.01605 (2024)

\bibitem{tihanyi2024cybermetric}
Tihanyi, N., Ferrag, M.A., Jain, R., Bisztray, T., Debbah, M.: CyberMetric: A benchmark dataset based on retrieval-augmented generation for evaluating LLMs in cybersecurity knowledge. In: 2024 IEEE International Conference on Cyber Security and Resilience (CSR) (2024)

\bibitem{alam2024ctibench}
Alam, M.T., Bhusal, D., Nguyen, L., Rastogi, N.: CTIBench: A benchmark for evaluating LLMs in cyber threat intelligence. In: Advances in Neural Information Processing Systems 37 (2024)

\bibitem{shao2024nyuctf}
Shao, M., Jancheska, S., Udeshi, M., Dolan-Gavitt, B., Xi, H., Milner, K., Chen, B., Yin, M., et al.: NYU CTF Bench: A scalable open-source benchmark dataset for evaluating LLMs in offensive security. In: Advances in Neural Information Processing Systems 37 (2024)

\bibitem{zhang2024cybench}
Zhang, A.K., Perry, N., Dulepet, R., Ji, J., Menders, C., Lin, J.W., Jones, E., Hussein, G., et al.: Cybench: A framework for evaluating cybersecurity capabilities and risks of language models. arXiv:2408.08926 (2024)

\bibitem{jing2024secbench}
Jing, P., Tang, M., Shi, X., Zheng, X., Nie, S., Wu, S., Yang, Y., Luo, X.: SecBench: A comprehensive multi-dimensional benchmarking dataset for LLMs in cybersecurity. arXiv:2412.20787 (2024)

\bibitem{zhang2024llmcyber}
Zhang, J., Bu, H., Wen, H., Liu, Y., Fei, H., Xi, R., Li, L., Yang, Y., et al.: When LLMs meet cybersecurity: A systematic literature review. arXiv:2405.03644 (2024)

\bibitem{pearce2022copilot}
Pearce, H., Ahmad, B., Tan, B., Dolan-Gavitt, B., Karri, R.: Asleep at the keyboard? Assessing the security of GitHub Copilot's code contributions. In: 2022 IEEE Symposium on Security and Privacy (SP) (2022)

\bibitem{perry2023insecure}
Perry, N., Srivastava, M., Kumar, D., Boneh, D.: Do users write more insecure code with AI assistants? In: Proceedings of the 2023 ACM SIGSAC Conference on Computer and Communications Security (2023)

\bibitem{deng2023pentestgpt}
Deng, G., Liu, Y., Mayoral-Vilches, V., Liu, P., Li, Y., Xu, Y., Zhang, T., Liu, Y., et al.: PentestGPT: An LLM-empowered automatic penetration testing tool. arXiv:2308.06782 (2023)

\bibitem{happe2023pwned}
Happe, A., Cito, J.: Getting pwn'd by AI: Penetration testing with large language models. In: Proceedings of the 31st ACM Joint European Software Engineering Conference and Symposium on the Foundations of Software Engineering (2023)

\bibitem{li2024wmdp}
Li, N., Pan, A., Gopal, A., Yue, S., Berrios, D., Gatti, A., Li, J.D., Dombrowski, A.-K., et al.: The WMDP benchmark: Measuring and reducing malicious use with unlearning. arXiv:2403.03218 (2024)

\bibitem{mazeika2024harmbench}
Mazeika, M., Phan, L., Yin, X., Zou, A., Wang, Z., Mu, N., Sakhaee, E., Li, N., et al.: HarmBench: A standardized evaluation framework for automated red teaming and robust refusal. arXiv:2402.04249 (2024)

\bibitem{yuan2024rjudge}
Yuan, T., He, Z., Dong, L., Wang, Y., Zhao, R., Xia, T., Xu, L., Zhou, B., et al.: R-Judge: Benchmarking safety risk awareness for LLM agents. Findings of the Association for Computational Linguistics: EMNLP 2024 (2024)

\bibitem{jiang2024wildteaming}
Jiang, L., Rao, K., Han, S., Ettinger, A., Brahman, F., Kumar, S., Mireshghallah, N., Lu, X., et al.: WildTeaming at scale: From in-the-wild jailbreaks to (adversarially) safer language models. In: Advances in Neural Information Processing Systems 37 (2024)

\end{thebibliography}

\end{document}
